% !TEX root = ../main.tex
\chapter{Nonlinear Modeling and Derivation}

\section{Kinematics}

With the pivot taken as the origin of an inertial frame, the ball center coordinates are

\begin{align}
  x_I &= r \cos \theta \\
  y_I &= r \sin \theta
\end{align}

where $r$ is the distance of the ball along the beam and $\theta$ is the beam angle from
horizontal.

Differentiating gives the inertial velocity components:

\begin{align}
  \dot{x}_I &= \dot{r}\cos\theta - r\dot{\theta}\sin\theta \\
  \dot{y}_I &= \dot{r}\sin\theta + r\dot{\theta}\cos\theta
\end{align}

The squared speed simplifies to the standard polar-coordinate result:

\begin{align}
  v^2
  &= \dot{x}_I^2 + \dot{y}_I^2 \\
  &= \left(\dot{r}\cos\theta - r\dot{\theta}\sin\theta\right)^2
   + \left(\dot{r}\sin\theta + r\dot{\theta}\cos\theta\right)^2 \\
  &= \dot{r}^2 + r^2 \dot{\theta}^2
\end{align}

The ball is modeled as rolling without slip along the beam, so its spin rate is

\begin{equation}
  \omega_{\text{ball}} = \frac{\dot{r}}{R}
\end{equation}

This removes the ball spin as an independent degree of freedom.

\section{Energy Terms}

\subsection{Ball Translational and Rotational Kinetic Energy}

The translational kinetic energy is

\begin{equation}
  T_{\text{trans}} = \frac{1}{2} m \left(\dot{r}^2 + r^2\dot{\theta}^2\right)
\end{equation}

Using the thin-shell inertia $I = \tfrac{2}{3}mR^2$ and the rolling constraint, the
rotational kinetic energy becomes

\begin{align}
  T_{\text{rot}}
  &= \frac{1}{2} I \omega_{\text{ball}}^2 \\
  &= \frac{1}{2}\left(\frac{2}{3}mR^2\right)\left(\frac{\dot{r}}{R}\right)^2 \\
  &= \frac{1}{3} m \dot{r}^2
\end{align}

\subsection{Beam Kinetic Energy}

The beam assembly rotates as a rigid body about the pivot:

\begin{equation}
  T_{\text{beam}} = \frac{1}{2} J \dot{\theta}^2
\end{equation}

\subsection{Total Kinetic Energy}

Summing all contributions gives

\begin{align}
  T
  &= T_{\text{trans}} + T_{\text{rot}} + T_{\text{beam}} \\
  &= \frac{5}{6}m\dot{r}^2 + \frac{1}{2}\left(mr^2 + J\right)\dot{\theta}^2
\end{align}

The $\tfrac{5}{6}$ factor on $\dot{r}^2$ is the characteristic effective-mass term of
the ball and beam plant.

\subsection{Potential Energy}

Taking the pivot height as the reference, the ball and beam potential energies are

\begin{align}
  V_{\text{ball}} &= mgr\sin\theta \\
  V_{\text{beam}} &= M_b g L_{\text{com}} \sin\theta
\end{align}

so the total potential energy is

\begin{equation}
  V = (mr + M_bL_{\text{com}}) g \sin\theta
\end{equation}

\subsection{Lagrangian}

The Lagrangian is therefore

\begin{equation}
  \mathcal{L}
  = \frac{5}{6}m\dot{r}^2
  + \frac{1}{2}\left(mr^2 + J\right)\dot{\theta}^2
  - (mr + M_bL_{\text{com}}) g \sin\theta
\end{equation}

\section{Euler--Lagrange Equations}

The generalized forces include viscous rolling damping and actuator torque:

\begin{align}
  Q_r &= -b_x \dot{r} \\
  Q_{\theta} &= \tau - b_{\theta}\dot{\theta}
\end{align}

\subsection{Ball Equation}

For the $r$ coordinate,

\begin{align}
  \frac{\partial \mathcal{L}}{\partial \dot{r}} &= \frac{5}{3}m\dot{r} \\
  \frac{d}{dt}\left(\frac{\partial \mathcal{L}}{\partial \dot{r}}\right) &= \frac{5}{3}m\ddot{r} \\
  \frac{\partial \mathcal{L}}{\partial r} &= mr\dot{\theta}^2 - mg\sin\theta
\end{align}

Applying the Euler--Lagrange equation yields

\begin{equation}
  \frac{5}{3}m\ddot{r} + b_x\dot{r} - mr\dot{\theta}^2 + mg\sin\theta = 0
  \label{eq:eom_r}
\end{equation}

The terms represent the effective rolling inertia, viscous damping, centrifugal
coupling from beam rotation, and gravitational acceleration along the beam.

\subsection{Beam Equation}

For the $\theta$ coordinate,

\begin{align}
  \frac{\partial \mathcal{L}}{\partial \dot{\theta}} &= (mr^2 + J)\dot{\theta} \\
  \frac{d}{dt}\left(\frac{\partial \mathcal{L}}{\partial \dot{\theta}}\right)
    &= (mr^2 + J)\ddot{\theta} + 2mr\dot{r}\dot{\theta} \\
  \frac{\partial \mathcal{L}}{\partial \theta}
    &= -(mr + M_bL_{\text{com}})g\cos\theta
\end{align}

Hence

\begin{equation}
  (J + mr^2)\ddot{\theta}
  + 2mr\dot{r}\dot{\theta}
  + b_{\theta}\dot{\theta}
  + (mr + M_bL_{\text{com}})g\cos\theta
  = \tau
  \label{eq:eom_theta}
\end{equation}

The beam equation captures inertia variation with ball position, Coriolis coupling,
pivot damping, gravity, and motor torque.

\section{Explicit Nonlinear State Model}

Defining the effective rolling mass

\begin{equation}
  m_e = \frac{5}{3}m = \SI{0.004667}{kg}
\end{equation}

the ball equation can be written explicitly as

\begin{equation}
  \ddot{r}
  = \frac{mr\dot{\theta}^2}{m_e}
  - \frac{mg}{m_e}\sin\theta
  - \frac{b_x}{m_e}\dot{r}
\end{equation}

and the beam equation becomes

\begin{equation}
  \ddot{\theta}
  =
  \frac{
    \tau
    - 2mr\dot{r}\dot{\theta}
    - b_{\theta}\dot{\theta}
    - (mr + M_bL_{\text{com}})g\cos\theta
  }{J + mr^2}
\end{equation}

Together, these equations form the coupled nonlinear plant with state
$[r,\dot{r},\theta,\dot{\theta}]^{\mathsf{T}}$. They are the source model used for the
linearized design discussion and for interpreting the runtime controller structure.
